\section{Complete checked proofs and source bridges}\label{app:lean}
The single research module imports only \texttt{Std}. The proof toolchain is
Lean 4.36.0-rc1, commit \texttt{c9afcddf514a6772bd4d649f7cfc04eff0b9f15e}.
Every finite statement below is proved by kernel reduction with
\texttt{decide}, not by a native decision oracle. The quantifier
\texttt{e : BitVec 7} covers every one of the $128$ subsets, including the
empty and full supports. The clean build and axiom reports are included
in the artifact. The permitted axioms are \texttt{propext},
\texttt{Classical.choice} and \texttt{Quot.sound}; no admitted proof is used.

\subsection{Public sources and the proof map}
The \href{https://github.com/alectionary/vibe-math/blob/master/formalization/gallager-b-absorbing/Counterexample.lean}{complete current Lean source}
and its
\href{https://github.com/alectionary/vibe-math/blob/2717b0ab2b59c2ccb71f0e369a9f21d80490c191/formalization/gallager-b-absorbing/Counterexample.lean}{immutable checked revision}
contain all research proofs, including the actual fourteen-coordinate bridge.
The
\href{https://github.com/alectionary/vibe-math/blob/master/formalization/gallager-b-absorbing/verification.json}{verification manifest},
\href{https://github.com/alectionary/vibe-math/blob/master/formalization/gallager-b-absorbing/prior-art.json}{source and search ledger}, and
\href{https://github.com/alectionary/vibe-math/blob/master/formalization/gallager-b-absorbing/verify.py}{independent literal-matrix verifier}
make the input domain, access limits and evidence distinctions explicit.

\small
\begin{longtable}{p{.27\textwidth}p{.61\textwidth}}
\toprule Manuscript proof & Matching Lean declarations (all in the linked module)\\
\midrule\endhead
Absorption and connectivity & \texttt{degrees}, \texttt{full\_set\_absorbing}, \texttt{first\_block\_connected}\\
Syndrome-column identity and Lemma~\ref{lem:votes} & \texttt{syndrome\_index\_bridge}, \texttt{simplex\_vote\_counts}, \texttt{first\_hard\_formula}, \texttt{first\_hard\_matches}\\
Extrinsic update, Lemma~\ref{lem:edges} & \texttt{first\_edge\_matches}, with the full \texttt{checkUpdate} and \texttt{variableUpdate} definitions\\
Theorem~\ref{thm:main}, eventual properties & \texttt{terminal\_matches}, \texttt{all\_eventually\_correct}, \texttt{all\_eventually\_converge}, \texttt{absorbing\_not\_trapping}\\
Steane miscorrection comparison & \texttt{steane\_logical\_miscorrection}\\
Lemma~\ref{lem:kernel} & \texttt{kernel\_equals\_four\_generator\_span}, \texttt{actual\_css\_commutation}\\
Proposition~\ref{prop:css}, complete span and nontrivial logical & \texttt{first\_block\_residual\_is\_stabilizer}, \texttt{first\_block\_quantum\_recovery}, \texttt{logical\_commutes}, \texttt{logical\_not\_stabilizer}\\
Actual source matrices and fourteen-coordinate update & \texttt{actual\_matrix\_first\_step\_bridge}, \texttt{actual\_matrix\_syndrome\_matches}, \texttt{terminal\_eq\_firstHard}, \texttt{actual\_matrix\_quantum\_recovery}\\
\bottomrule
\end{longtable}
\normalsize

\lstset{escapeinside={(*@}{@*)}}

\subsection{Matrix, subsets and exact update rules}
A seven-bit word is the indicator of a subset of the seven actual columns.
The matrix is the binary dot product on labels $1,\ldots,7$.
\texttt{pack\_bits} checks that packing and reading these column bits agree.
The two message functions exclude the recipient explicitly, and the
majority definition implements zero ties by a strict inequality.
Here are the complete definitions and finite source checks. Unicode symbols are rendered typographically; the linked Lean file is the checked source.
\lstinputlisting[firstline=1,lastline=81]{proof-listing.txt}

\subsection{From accepted output to absence of trapping}
The dependency chain is
\texttt{first\_hard\_matches} $\to$ \texttt{terminal\_matches} $\to$
the two eventual properties $\to$ \texttt{absorbing\_not\_trapping}.
The output at time zero is the initial zero guess. From time one onward
it is the accepted terminal value. Thus the following unbounded-in-time
proof is a structural argument, rather than a finite simulation horizon.
The complete central proof bodies are printed here.
\lstinputlisting[firstline=83,lastline=126]{proof-listing.txt}
Direct declaration links:
\href{https://github.com/alectionary/vibe-math/blob/2717b0ab2b59c2ccb71f0e369a9f21d80490c191/formalization/gallager-b-absorbing/Counterexample.lean#L101}{eventual check correctness},
\href{https://github.com/alectionary/vibe-math/blob/2717b0ab2b59c2ccb71f0e369a9f21d80490c191/formalization/gallager-b-absorbing/Counterexample.lean#L108}{eventual variable convergence}, and
\href{https://github.com/alectionary/vibe-math/blob/2717b0ab2b59c2ccb71f0e369a9f21d80490c191/formalization/gallager-b-absorbing/Counterexample.lean#L115}{nontrapping theorem}.

\subsection{Full CSS span and logical witness}
A pair of seven-bit words is exactly a binary word on the disjoint union
of two seven-coordinate blocks. The matrices below retain all fourteen
$X$ checks and all seven $Z$ generators. \texttt{combine} enumerates the
entire generator span with independent binary coefficients; it does not
check only individual stabilizer rows. The recovery statement covers every
input and witnesses membership of the residual in that whole span.
The nontrivial logical test uses the same matrix and span definitions.
\lstinputlisting[firstline=128,lastline=191]{proof-listing.txt}
Direct declaration links:
\href{https://github.com/alectionary/vibe-math/blob/2717b0ab2b59c2ccb71f0e369a9f21d80490c191/formalization/gallager-b-absorbing/Counterexample.lean#L151}{complete residual-span theorem} and
\href{https://github.com/alectionary/vibe-math/blob/2717b0ab2b59c2ccb71f0e369a9f21d80490c191/formalization/gallager-b-absorbing/Counterexample.lean#L165}{nontrivial logical theorem}.

\subsection{Analytic identities and the actual fourteen-coordinate decoder}
The finite formulas below supply the syndrome-column, vote-count and
kernel/span identities used in the written proofs. The final portion then
calculates the syndrome and all first-step votes directly from the actual
fourteen-coordinate rows. Its bridge theorem proves that this calculation
is the component calculation used above. This avoids assuming that an
abstract block argument automatically matches the source matrix.
\lstinputlisting[firstline=193,lastline=274]{proof-listing.txt}
Direct declaration link:
\href{https://github.com/alectionary/vibe-math/blob/2717b0ab2b59c2ccb71f0e369a9f21d80490c191/formalization/gallager-b-absorbing/Counterexample.lean#L259}{actual-matrix quantum recovery theorem}.

\subsection{Reproduction and evidence limits}
Run \texttt{python3 build.py --lean /path/to/lean} from the artifact
folder with the pinned toolchain, then \texttt{python3 verify.py --check}.
The Lean result is the proof; the independently regenerated JSON certificate
provides a second, literal-matrix reconstruction of all inputs.
No sampled experiment is used. The certificate also supplies independent
linear-algebra diagnostics; numerical diagnostics not stated as theorems
in the manuscript are not additional formal results.
The source ledger distinguishes the accessible arXiv source, the restricted
journal text, partial classical-source retrievals and historical questions
harvested without claiming they remain open today.
